\section{Method}
\label{sec:method}

\subsection{Dataset}
\label{sec:data}

Data were collected from keybr.com, an adaptive typing practice platform,
by a single participant over 17 months (December 2024 -- May 2026).
Each lesson record contains: UTC timestamp, overall typing speed (characters
per minute), error count, session duration (ms), session length (characters),
and per-key average keystroke time and error count for 27 characters
(a--z and space).
After timestamp parsing and WPM conversion (\(\text{WPM} = \text{cpm}/5\)),
the dataset contains \textbf{42,501 lessons} with a mean speed of 87.6 WPM
and a peak of 136.6 WPM.
No personally identifiable information beyond typing behavior is present.

\subsection{Feature Engineering}
\label{sec:features}

We extract 12 behavioral features organized into three groups.

\paragraph{Speed and Variance Features.}
\emph{Residual Execution Velocity} (\textbf{EV}) captures per-lesson speed
surprise relative to recent trend:
\begin{equation}
  \text{EV}_t = \text{WPM}_t - \frac{1}{7}\sum_{i=0}^{6}\text{WPM}_{t-i}
  \label{eq:ev}
\end{equation}
Positive EV indicates a faster-than-trend session; negative indicates drift.
\emph{Rolling Stability Variance} (\textbf{RSV}) is the 30-lesson rolling
standard deviation of WPM, measuring consistency over time.
\emph{WPM vs.\ 50-lesson rolling} (\textbf{wpm\_vs\_50roll}) normalizes
current WPM against the long-term trend.

\paragraph{Error Features.}
\emph{Residual Error-Speed Recovery Penalty} (\textbf{ESRP}) captures
per-lesson error surprise:
\begin{equation}
  \text{ESRP}_t = r_t \cdot \frac{1000}{\text{WPM}_t} -
  \frac{1}{5}\sum_{i=0}^{4}\!\left(r_{t-i}\cdot\frac{1000}{\text{WPM}_{t-i}}\right)
  \label{eq:esrp}
\end{equation}
where \(r_t = \text{errors}_t / \text{chars}_t\) is the per-lesson error rate.
\emph{Error Entropy} (\textbf{error\_entropy}) is the Shannon entropy of the
per-key error count distribution, measuring whether errors are concentrated on
specific keys or spread uniformly.

\paragraph{Structural Features.}
\emph{Sequential Stamina Fatigue} (\textbf{SSF}) encodes intra-day fatigue as
the product of the lesson's position within the day and the within-day
WPM z-score.
\emph{Circadian Temporal Anchors} (\textbf{CTA}, \textbf{CTA\_cos}) are the
sine and cosine of the session's UTC hour, encoding time-of-day in a
continuous circular representation.
\emph{Lesson Duration} (\textbf{lesson\_dur\_z}) is \(\log_{1\!+}(session\_time)\),
log-transforming the heavy-tailed duration distribution.
\emph{Active Character Fraction} (\textbf{active\_char\_frac}) is the
proportion of the 27-character vocabulary present in the lesson, encoding
curriculum depth.
\emph{Hand Asymmetry} (\textbf{hand\_asymmetry}) is the ratio of mean
left-hand to right-hand key latency across available keys, capturing
motor lateralization.
\emph{Practice Gap} (\textbf{practice\_gap}) is \(\log_{1\!+}(h)\) where
\(h\) is the elapsed hours since the previous lesson (clipped to 72h),
encoding schedule regularity.

\paragraph{The Residual Feature Insight.}
\label{sec:residual}
Prior work defined EV as the 7-lesson rolling mean WPM.
This creates a \emph{structural reconstruction ceiling}: at inference time,
the model receives a single lesson's feature vector.
Reconstructing a rolling mean requires access to the 6 preceding lessons ---
information unavailable in the bottleneck.
The autoencoder therefore cannot learn to reconstruct rolling-mean EV from
single-lesson features, regardless of capacity or training time.

By reformulating as the \emph{deviation} from the rolling mean
(Eqs.~\ref{eq:ev}--\ref{eq:esrp}), EV and ESRP become single-lesson
quantities that \emph{are} predictable from other within-lesson features
(WPM magnitude, error rate, circadian time).
This yields a theoretically principled formulation that is empirically
validated in Section~\ref{sec:ablation}.

\subsection{Architecture}
\label{sec:arch}

We use a symmetric undercomplete autoencoder:
\[
12 \xrightarrow{} 32 \xrightarrow{} 16 \xrightarrow{} \mathbf{10}
  \xrightarrow{} 16 \xrightarrow{} 32 \xrightarrow{} 12
\]
with 2,230 trainable parameters.
The encoder uses LeakyReLU activations (\(\alpha{=}0.01\)); the 10-dimensional
bottleneck uses Tanh to bound the \emph{Muscle Memory Signature}.
The decoder mirrors the encoder with LeakyReLU activations and a linear
output layer.
All weights are initialized with Xavier uniform initialization.
The bottleneck compression ratio (10/12\,=\,0.83) was selected to maintain
sufficient capacity for 12 features while still forcing compression.

A \emph{reproducible denoising mask} applies 10\% random feature zeroing
during training using a deterministic random state seeded by epoch and
batch index, ensuring reproducibility under DP noise.

\subsection{Differentially Private Training}
\label{sec:dp}

We train with DP-SGD via the Opacus library~\cite{yousefpour2021opacus}
using the RDP accountant.
Hyperparameters: target \(\varepsilon{=}8.0\), \(\delta{=}10^{-5}\),
max gradient norm \(C{=}1.0\), batch size 64, yielding noise multiplier
\(\sigma{=}0.5994\).
The final privacy expenditure is \(\varepsilon{=}7.9983\) at \(\delta{=}10^{-5}\).

The weighted MSE loss emphasizes features with diagnostic importance:
\begin{equation}
  \mathcal{L} = \frac{1}{N}\sum_{i=1}^{N} \mathbf{w}^\top (\hat{\mathbf{x}}_i - \mathbf{x}_i)^2
  \label{eq:loss}
\end{equation}
where \(\mathbf{w} \in \mathbb{R}^{12}\) is a fixed weight vector.
Weights are set to 4.0 for lesson\_dur\_z (hardest to reconstruct),
2.5 for practice\_gap, 2.0 for RSV, hand\_asymmetry, and error\_entropy,
and 1.0 for EV and ESRP (residuals are structurally easy to predict).

We use the Adam optimizer with initial learning rate \(2\times10^{-3}\) and
\texttt{ReduceLROnPlateau} scheduling (factor\,=\,0.5, patience\,=\,5),
which halves the learning rate when validation loss stagnates.
This allows the model to escape early local minima under DP noise by
adaptively reducing gradient step size, in contrast to cosine annealing which
fixes the decay schedule.
Training uses a patience of 20 epochs for early stopping with best-checkpoint
restoration.

\subsection{Longitudinal Drift Index and Anomaly Detection}
\label{sec:ldi}

For each lesson, per-lesson reconstruction MSE is computed as:
\[
  \ell_t = \frac{1}{12}\sum_{j=1}^{12}(\hat{x}_{t,j} - x_{t,j})^2
\]
Daily LDI is the mean per-lesson MSE across all lessons in a calendar day,
smoothed with a 7-day rolling mean to suppress lesson-level noise:
\[
  \text{LDI}_{7\text{d}}(d) = \frac{1}{7}\sum_{k=0}^{6}\overline{\ell}_{d-k}
\]

\paragraph{Bootstrap Confidence Interval Alerting.}
Rather than a single point threshold, we estimate the uncertainty in the
training-period anomaly threshold.
We draw 2,000 bootstrap resamples of training-period LDI values and compute
the \(\mu + 2\sigma\) threshold for each resample.
The 2.5th and 97.5th percentiles of these thresholds define the lower and
upper CI bounds \((\theta_{\text{lo}}, \theta_{\text{hi}})\).
Each day is then classified as:
\[
\text{tier}(d) =
\begin{cases}
  \text{HIGH}     & \text{if } \text{LDI}_{7\text{d}}(d) \geq \theta_{\text{hi}} \\
  \text{MARGINAL} & \text{if } \theta_{\text{lo}} \leq \text{LDI}_{7\text{d}}(d) < \theta_{\text{hi}} \\
  \text{NORMAL}   & \text{if } \text{LDI}_{7\text{d}}(d) < \theta_{\text{lo}}
\end{cases}
\]
HIGH days represent confident anomalies; MARGINAL days fall within the
threshold uncertainty band and warrant monitoring but not immediate action.
